\documentclass[11pt,a4paper]{article}
\usepackage[utf8]{inputenc}
\usepackage[T1]{fontenc}
\usepackage[margin=1in]{geometry}
\usepackage{hyperref}
\usepackage{listings}
\usepackage{xcolor}
\usepackage{graphicx}
\usepackage{booktabs}
\usepackage{amsmath}
\usepackage{enumitem}

\hypersetup{
    colorlinks=true,
    linkcolor=blue,
    urlcolor=blue,
    citecolor=blue
}

\definecolor{codebg}{RGB}{245,245,245}
\lstset{
    basicstyle=\ttfamily\small,
    backgroundcolor=\color{codebg},
    frame=single,
    breaklines=true,
    numbers=left,
    numberstyle=\tiny
}

\title{\textbf{Chat App: Technical Overview} \\ \large A Real-Time Messaging System}
\author{}
\date{}

\begin{document}

\maketitle
\tableofcontents
\newpage

\section{Project Overview}

The Chat App is a production-grade real-time messaging application designed to demonstrate competencies in \textbf{distributed systems}, \textbf{concurrent programming}, and \textbf{full-stack development}. The system enables multiple users to communicate in real time within named chat rooms, with support for multiple concurrent rooms and graceful handling of client connections and disconnections.

The project comprises two primary implementations:
\begin{enumerate}
    \item A \textbf{modern stack} (Go + React) using WebSockets for real-time bidirectional communication
    \item A \textbf{systems-level} implementation in C using POSIX sockets and pthreads
\end{enumerate}

Both implementations share the same conceptual model: a central server that accepts client connections, maintains room state, and broadcasts messages to all participants within a room.

\section{Tech Stack}

\subsection{Backend: Go (Golang)}

\textbf{What:} Go 1.21 with the standard library and Gorilla WebSocket package.

\textbf{Why:}
\begin{itemize}
    \item \textbf{Concurrency model:} Go's goroutines and channels provide a natural fit for handling thousands of concurrent WebSocket connections. Each client runs in its own goroutine without the overhead of OS threads.
    \item \textbf{Simplicity:} Go's design favors readability and maintainability. The server logic is concise and easy to reason about.
    \item \textbf{Performance:} Compiled to native code with low memory footprint; well-suited for I/O-bound workloads like chat servers.
    \item \textbf{Cross-platform:} Single codebase compiles to Linux, Windows, and macOS without modification.
\end{itemize}

\subsection{WebSocket Protocol}

\textbf{What:} The WebSocket protocol (RFC 6455) provides full-duplex communication over a single TCP connection.

\textbf{Why:}
\begin{itemize}
    \item \textbf{Real-time:} Unlike HTTP request-response, WebSockets allow the server to push messages to clients instantly without polling.
    \item \textbf{Efficiency:} After the initial handshake, overhead is minimal (small binary frames). No repeated HTTP headers.
    \item \textbf{Browser support:} Native \texttt{WebSocket} API in all modern browsers; no additional libraries required on the client.
\end{itemize}

\subsection{Gorilla WebSocket}

\textbf{What:} A popular, well-maintained Go implementation of the WebSocket protocol (\url{https://github.com/gorilla/websocket}).

\textbf{Why:}
\begin{itemize}
    \item Reference implementation for Go; widely used in production.
    \item Handles the HTTP-to-WebSocket upgrade, frame parsing, and connection lifecycle.
    \item Provides \texttt{Upgrader} for upgrading HTTP connections and \texttt{Conn} for reading/writing messages.
\end{itemize}

\subsection{Frontend: React + TypeScript}

\textbf{What:} React 18 with TypeScript, built with Vite.

\textbf{Why:}
\begin{itemize}
    \item \textbf{React:} Component-based architecture; efficient DOM updates via virtual DOM. Hooks (\texttt{useState}, \texttt{useEffect}, \texttt{useCallback}) simplify state and side-effect management for the WebSocket connection.
    \item \textbf{TypeScript:} Static typing catches errors at compile time. Interfaces for \texttt{ConnectionConfig}, \texttt{WSMessage} document the data flow and improve maintainability.
    \item \textbf{Vite:} Fast development server with HMR (Hot Module Replacement). ESBuild-based builds for quick production bundles.
\end{itemize}

\subsection{Message Protocol: JSON}

\textbf{What:} Messages are serialized as JSON with a \texttt{type} field and optional \texttt{username}, \texttt{content}, \texttt{timestamp}, \texttt{room}.

\textbf{Why:}
\begin{itemize}
    \item Human-readable and easy to debug.
    \item Native support in JavaScript (\texttt{JSON.parse}, \texttt{JSON.stringify}) and Go (\texttt{encoding/json}).
    \item Extensible: new message types can be added without breaking existing clients.
\end{itemize}

\subsection{Deployment: Docker}

\textbf{What:} Docker Compose with multi-stage Dockerfiles for the Go server and React client (served via Nginx).

\textbf{Why:}
\begin{itemize}
    \item \textbf{Reproducibility:} One command (\texttt{docker compose up}) runs the entire stack. No ``works on my machine'' issues.
    \item \textbf{Multi-stage builds:} Builder stage compiles the application; runtime stage uses a minimal Alpine image. Smaller, more secure containers.
    \item \textbf{Nginx:} Serves static React build and proxies \texttt{/ws} to the Go server, so the client connects to a single origin (no CORS complexity in production).
\end{itemize}

\subsection{C Implementation (Systems Showcase)}

\textbf{What:} A separate implementation in C using POSIX sockets (\texttt{sys/socket.h}, \texttt{netinet/in.h}), pthreads, and mutexes.

\textbf{Why:}
\begin{itemize}
    \item \textbf{Systems understanding:} Demonstrates low-level socket programming, TCP stream semantics, and manual memory management.
    \item \textbf{Concurrency:} One thread per client; \texttt{pthread\_mutex\_lock} protects the shared \texttt{client\_sockets} array from race conditions.
    \item \textbf{Educational:} Contrasts with the Go implementation—same problem, different abstraction levels.
\end{itemize}

\section{Architecture}

\subsection{High-Level Flow}

\begin{enumerate}
    \item User enters username, room name, and server URL in the Lobby.
    \item Client opens WebSocket: \texttt{ws://host/ws?username=X\&room=Y}.
    \item Server upgrades HTTP to WebSocket, creates a \texttt{Client}, registers it with the \texttt{Hub}.
    \item Hub adds client to the room, broadcasts ``X joined'' to existing members, sends member list to the new client.
    \item User types a message; client sends \texttt{\{"type":"chat","content":"..."\}}.
    \item Hub receives via \texttt{broadcast} channel, forwards to all clients in the room except the sender.
    \item On disconnect, Hub unregisters client, broadcasts ``X left'', cleans up.
\end{enumerate}

\subsection{Hub Pattern}

The Hub is the central coordinator:
\begin{itemize}
    \item \textbf{rooms:} \texttt{map[string]map[string]*Client} — room name $\to$ client ID $\to$ client.
    \item \textbf{Channels:} \texttt{register}, \textbf{unregister}, \texttt{broadcast} — all state changes go through these channels, ensuring a single goroutine modifies shared data (no locks needed for the hub loop).
    \item \textbf{Mutex:} \texttt{sync.RWMutex} protects the \texttt{rooms} map during iteration and lookup from the channel handlers.
\end{itemize}

\subsection{Project Structure}

\begin{verbatim}
Chat App/
├── server/                 # Go WebSocket server
│   ├── main.go             # HTTP server, /ws route
│   └── internal/
│       ├── hub/            # Room & client management
│       └── websocket/      # Connection upgrade & handler
├── client/                 # React + TypeScript
│   ├── src/
│   │   ├── components/     # Lobby, ChatRoom
│   │   ├── hooks/         # useWebSocket
│   │   └── App.tsx
│   └── package.json
├── c-server/               # C implementation
│   ├── server.c
│   ├── client.c
│   └── Makefile
├── docker-compose.yml
├── Dockerfile.server
└── Dockerfile.client
\end{verbatim}

\section{Key Design Decisions}

\begin{table}[h]
\centering
\begin{tabular}{llll}
\toprule
\textbf{Decision} & \textbf{Choice} & \textbf{Alternative} & \textbf{Reason} \\
\midrule
Transport & WebSocket & HTTP long-polling & Lower latency, less overhead \\
Backend language & Go & Node.js, Rust & Concurrency, simplicity, deployability \\
State coordination & Channels & Shared memory + locks & Avoids deadlocks, clearer flow \\
Message format & JSON & Protocol Buffers & Simplicity, browser support \\
Frontend & React + TS & Vue, Svelte & Ecosystem, type safety \\
\bottomrule
\end{tabular}
\caption{Design decisions and rationale}
\end{table}

\section{Conclusion}

The Chat App demonstrates a complete real-time messaging system from low-level sockets (C) to high-level abstractions (Go, React). The tech stack choices prioritize \textbf{simplicity}, \textbf{performance}, and \textbf{maintainability}, while the dual implementation (Go + C) showcases both software engineering and systems engineering competencies.

\end{document}
